\begin{figure}[tbp]
\centering
\begin{tikzpicture}
\begin{axis}[
  width=0.82\textwidth,height=8.4cm,
  xmin=-1.8,xmax=0.6,ymin=-1.5,ymax=1.5,
  xlabel={$\log([\mathrm{N\,II}]\,6583/\mathrm{H}\alpha)$},
  ylabel={$\log([\mathrm{O\,III}]\,5007/\mathrm{H}\beta)$},
  grid=major,grid style={AtlasNavy!10},
  axis line style={AtlasNavy},
  label style={font=\small},tick label style={font=\small}]
  \addplot[AtlasCyan,line width=1.5pt,domain=-1.75:-0.08,samples=150,unbounded coords=jump] {0.61/(x-0.05)+1.3};
  \addplot[AtlasAmber,line width=1.5pt,dashed,domain=-1.75:0.30,samples=150,unbounded coords=jump] {0.61/(x-0.47)+1.19};
  \node[AtlasNavy] at (axis cs:-1.15,-0.65){star forming};
  \node[AtlasNavy] at (axis cs:0.12,0.95){AGN-like};
  \node[AtlasNavy,rotate=57] at (axis cs:-0.45,0.15){composite};
\end{axis}
\end{tikzpicture}
\caption{The [N~II] BPT diagram. The solid empirical curve follows \citet{kauffmann2003}; the dashed theoretical ``maximum starburst'' curve follows \citet{kewley2001}. Classification must account for uncertainties, mixing, diffuse ionized gas, shocks, and aperture.}
\label{fig:bpt}
\end{figure}
